% 384_FFGFT_Kurzfassung_En_ch.tex
% Johann Pascher, 1 Oct. 2026
% FFGFT in brief: foundations, results, status (after the calculation review)
% Revised 1 Oct. 2026: emphasis on ratios and natural units

\providecommand{\BM}{\textbf{[B]}}
\providecommand{\KM}{\textbf{[K]}}
\providecommand{\SM}{\textbf{[S]}}
\providecommand{\XM}{\textbf{[X]}}
\providecommand{\QM}{\textbf{[Q]}}

\chapter*{FFGFT in Brief}
\addcontentsline{toc}{chapter}{FFGFT in Brief}

\begin{abstract}
The Fundamental Fractal Geometric Field Theory (FFGFT) is first and foremost a ratio-based
theory. It works in natural units, $\hbar=c=1$, with the convention $\alpha=1$ and a
correspondingly redefined charge $e=\sqrt{4\pi}$. In this basic form there is exactly one
dimensionless parameter, $\xi=4/30000=1/7500$, and no fractal correction. Particles, couplings
and gravitation appear as properties of a compact four-dimensional carrier $T^4/\mathbb{Z}_3$;
the starting point is the time-mass duality $\tilde T\cdot m=1$. The core statements are ratios:
$m_\mu/m_e$ and $m_\tau/m_\mu$ follow from $\xi$ to within $0.5$ and $1\,\%$, the electroweak
scale relates to the Planck energy as $\xi^4/(5\pi)$, the minimal length and minimal time relate
to the Planck quantities as $\xi$, the neutrino mass differences agree to $1\,\%$, and charge
quantisation follows algebraically from the Galois structure of $\mathrm{GF}(27)^*$. Only the
translation into SI needs a measured value as anchor, the factor $K_{\text{frak}}=74/75$ and
conversion quantities such as $C_{\text{conv}}$; this is also where the absolute values of
$\alpha$ and $G$ arise. This document summarises the corpus (Docs.~001--386, A series) as of
1~October 2026, after the complete recalculation. For every result it gives the formula, the
numerical value, the comparison value, the status and the documents that carry it, and it states
just as clearly what is set by choice, what remains open and what no longer holds. The cosmic
sector ($H_0$, $\Lambda$) is deliberately set aside.
\end{abstract}

% ============================================================
\section*{Status markers}
% ============================================================
\begin{center}
\begin{tabular}{@{}lp{11cm}@{}}
\textbf{[B]} & Algebraically proved (transformation or proof without additional assumption). \\[2pt]
\textbf{[K]} & Numerically confirmed (verification script, comparison with measurements). \\[2pt]
\textbf{[S]} & Choice, open or speculative: motivated but not derived. \\[2pt]
\textbf{[X]} & Refuted or excluded. \\[2pt]
\textbf{[Q]} & Convention or adopted quantity. \\
\end{tabular}
\end{center}
The corrected state is authoritative: the dated notes in the documents (``Corrected on'',
``Note of'', ``Removed on'') and the correction register Doc.~190 up to R146. Comparison values:
CODATA 2022, PDG 2024. All numbers in this document are recomputed by the verification script
\texttt{pruef\_384\_zusammenfassung.py}.

% ============================================================
\section{What FFGFT is}
% ============================================================

\subsection{Ratios in natural units}

The core of FFGFT is ratios. In natural units, mass, energy, temperature and inverse time are one
quantity in four guises (A085), and every physical statement can be written as a dimensionless
ratio: mass ratios, scale ratios to the Planck quantities, mixing angles, charge ratios. Such
ratios do not depend on any choice of units and carry the content of the theory. Absolute values
in SI are a second layer: they need a measured anchor and conversion quantities, and corrections
appear there that do not occur in the ratios (A080, A130, Doc.~122). This summary follows that
order: first the basic form and the ratios, then the translation into SI.

\subsection{Architecture in four layers}

FFGFT rests on three explicitly chosen basic assumptions (R56, Doc.~365): the time-mass duality
$\tilde T\cdot m=1$, the compact carrier $T^4$ and the identification with order~3,
$\mathbb{Z}_3$. A fourth, motivated structural choice is the lattice $D_4$; Doc.~380 shows that
among the comparison carriers tested with a $\mathbb{Z}_3$ action that has fixed points, $D_4$ has
the lowest lattice energy \KM, but the choice remains a choice \SM. In addition there is exactly
one dimensionless parameter, $\xi=4/30000$. The fine-structure constant is set to $\alpha=1$ by
adjusting the charge unit; this is a convention, not a result \QM. Galois fields, Hodge theory
and representation theory serve as consistency checks, not as additional inputs (Doc.~365,
layer~4). The spacetime dimension $3+1$ is an empirical input (P34).

\subsection{Guidelines}

Four rules determine how results are to be read. First, ratios are what counts; an absolute value
in SI tests the conversion, not the structure (A130). Second, FFGFT works with $\xi$ alone;
measured values from SI and the Standard Model serve only for conversion and comparison, and for
that exactly one measured value is needed as anchor (Doc.~383). Third, integer structural choices
such as $\mathbb{Z}_3$, $D_4$, the winding numbers or the exponent~10 in the cosmic sector are
motivated choices, not calibrations. Fourth, the cosmic sector ($H_0$, $\Lambda$) is deliberately
excluded from the derivation of constants (P39): the Standard Model does not derive it either, and
none of the frameworks compared determines the ratio of Planck length to Hubble length from first
principles.

\subsection{What the corpus no longer claims}

The calculation review has withdrawn several earlier claims. FFGFT is not ``postulate-free'':
$\tilde T\cdot m=1$, $T^4$, $\mathbb{Z}_3$, $D_4$ and $\xi$ are chosen. Precision tables of
older documents in which prefactors had been fitted to measured values are not evidence (R138).
The agreement of the SI value of $\alpha$ is a property of the chosen anchor, not an independent
prediction (R135). The precision of $G$ to $0.003\,\%$ is not evidence in its own right (R141). A
$\xi$ effect in Bell experiments or quantum gates that is measurable with present-day hardware
does not exist (R142, R143). Where such statements appear in older documents, they carry a dated
note.

% ============================================================
\section{The basic form}
% ============================================================

\subsection{Natural units, \texorpdfstring{$\alpha=1$}{alpha = 1} and the charge}

With $\hbar=c=\varepsilon_0=1$ (Heaviside--Lorentz) what remains of the fine-structure constant
is the relation $e^2=4\pi\alpha$ \BM. The basic form sets $\alpha=1$ \QM\ and hence
$e=\sqrt{4\pi}$: the charge is measured by the full surface of the sphere (A267). $\alpha$ does
not disappear in the process; it moves into the charge unit,
\begin{equation}
  \alpha=\Bigl(\frac{e_{\text{hist}}}{e_{\text{geo}}}\Bigr)^2=\frac{r_e}{\lambda_C}\quad\BM ,
\end{equation}
the ratio of the Coulomb and Compton spheres of the electron (Doc.~386). Geometrically,
$\alpha=1$ means that both spheres coincide, and with them the Bohr radius. In every calculation
$\alpha$ stays on record through $e^2=4\pi\alpha$, one factor $\sqrt\alpha$ per vertex (A267). The
bridge between $\alpha$ and $\xi$ runs through the energy scale $E_0=\sqrt{m_em_\mu}$, the
geometric mean of the two lightest charged leptons; in the basic form it reads
\begin{equation}
  \xi\,E_0^2=1,\qquad E_0^2=\frac1\xi=7500\quad\BM .
\end{equation}
This fixes the T0 energy unit and leaves no room for a correction factor.

\subsection{Time-mass duality}

In natural units the intrinsic time scale and the mass of a state are reciprocal:
\begin{equation}
  \tilde T\cdot m=1,\qquad \tilde T=\frac{\hbar}{mc^2}\ \text{(SI)}.
\end{equation}
The relation describes every state as a snapshot, not as a law of motion (Doc.~365). A
position-dependent mass in a gravitational potential is equivalent to a position-dependent time;
``mass constant, space curved'' and ``space flat, mass variable'' are two representations of the
same situation (Doc.~078 \KM). On the carrier the duality holds mode by mode exactly as a position
relation $\lambda_4\cdot m=2\pi$ (Doc.~314 \KM). Reciprocity gives
$E\to\infty\Leftrightarrow T\to0$; since time is bounded from below by $T_0=\xi\,t_P$, the
horizon is an extremum and not a singularity (Doc.~306).

\subsection{The parameter \texorpdfstring{$\xi$}{xi}}

\begin{equation}
  \xi=\frac{4}{30000}=\frac{1}{7500}=1.\overline{3}\times10^{-4}.
\end{equation}
$\xi$ is the only dimensionless parameter. It has several consistent readings, but none produces
the number without an input: geometrically as the number~4 of lattice directions of $D_4$ over
$3\cdot10^4$ (Docs.~004, 320); as a Galois number via
\begin{equation}
  \frac{(r_\mu/r_e)^2}{\xi}=|\mathrm{GF}(9)^*|^2\cdot5^2\cdot|\mathrm{GF}(27)|=8^2\cdot25\cdot27=43200
\end{equation}
(Docs.~336, 338), which is arithmetically exact \BM\ but depends on the chosen prefactors $r_i$;
and as the smallest eigenvalue of the $D_4$ sub-operator, $\xi=\lambda_{\min}(\hat F_{D_4})$
(Docs.~322, 327), which is well defined provided the $D_4$ sector lies in the $\mathbb Z_3$-invariant
sector \BM, but whose value is inserted, not derived. Overall, $\xi$ is a motivated choice \SM: in
the words of Doc.~365, not completely free, but not yet completely derived either.

The Higgs geometry serves as a check, and it too is a pure ratio: with the self-coupling
$\lambda_h=m_h^2/(2v^2)=0.129$ one has
$\xi_{\text{EFT}}=\lambda_h^2v^2/(16\pi^3m_h^2)=m_h^2/(64\pi^3v^2)=1.30\times10^{-4}$; the
expression depends only on $m_h/v$ and lies about $2.3\,\%$ below $4/30000$ (PDG 2024;
$2.2$--$2.6\,\%$ depending on rounding; Docs.~354, A190, 385) \KM; the deviation is interpreted
structurally \SM\ (R59). This is a consistency check to within a few per cent, not an exact
derivation. The route goes back to the original vacuum formula of 2025 with $\varepsilon_0$;
written out it is $\xi_{\text{EFT}}\cdot\alpha$, and with $\alpha=1$ the same expression
(Doc.~385) \BM. Only the shortcut with $\lambda_h=m_h/v$ to $1/(16\pi^3)$, 15 times $\xi$
(Doc.~320), does not hold \XM.

\subsection{Scales as ratios}

All scales of the theory stand in fixed ratios to one another. The minimal length, minimal time
and maximal energy relate to the Planck quantities as
\begin{equation}
  \frac{L_0}{\ell_P}=\frac{T_0}{t_P}=\xi,\qquad \frac{E_{\max}}{E_P}=\frac1\xi=7500 .
\end{equation}
The exponent~1 is derived \BM\ (R90): $L_0$ is the Compton range of the time-field mediator with
mass $m_T=M_{\text{basis}}/\xi$; the choice $M_{\text{basis}}=m_P$ in the fundamental sector is
open \SM. $L_0$ makes the number of levels of the fractal operator finite and hence $\hat F$
bounded \BM\ (Doc.~327). The electroweak scale relates to the Planck energy, and the lepton masses
via the ladder (Section~3), likewise in pure $\xi$ ratios (Docs.~149, 383):
\begin{equation}
  \frac{v}{E_P}=\frac{\xi^4}{5\pi}=2.012\times10^{-17}\ (-0.23\,\%),\qquad
  \frac{m_i}{E_P}=\frac{r_i}{5\pi}\,\xi^{p_i+4}\quad\BM .
\end{equation}
With the measured $E_P$ the lepton masses come out as $0.504$, $104.8$ and $1780.9$\,MeV
($-1.32$, $-0.80$, $+0.22\,\%$) \KM. Open is the factor~10 in
$v=E_P/(f^4\cdot\tfrac{\pi}{2}\cdot10)$ with $f=1/\xi$; it is motivated but not derived \SM.

\subsection{Time structure and recursion}

The time structure runs step by step:
\begin{equation}
  \xi_{n+1}=\xi_n\,(1-100\,\xi_n),\qquad \xi_0=\tfrac{1}{7500},\qquad
  \xi_n\approx\frac{1}{100(n+75)},\qquad D(k)=\sum_{n<k}100\,\xi_n\approx\ln\!\Bigl(1+\frac{k}{75}\Bigr).
\end{equation}
In the frozen case with constant defect $d=1/75$ the winding closes after exactly 75 turns; the
rotation number is $74/75=1-100\xi$ \BM. This is the geometric origin of the value that appears
as $K_{\text{frak}}$ in the translation into SI (Section~4). In the running case the time winding
becomes a logarithmic spiral with self-similarity ratio $e$, a continuous scale invariance with a
pole at $n=-75$ (Doc.~295 \BM/\KM). The increments $d_n=100\,\xi_n$ decrease monotonically from
$d_1=1/75$; $d_1$ is the largest, not the smallest increment, and the lower time bound is
$T_0=\xi t_P$ (Docs.~306, 307). Doc.~381 shows the telescoping product
$\prod(1-100\xi_k)=\xi_n/\xi_0$ exactly \BM, a second order
$1/\xi_n\approx100(n+75)+100\ln((n+75)/75)$ that is 50 to about 1000 times more accurate \KM, and
the constant $\psi'(75)=0.01342$ between mass and time defect \KM. The assignment to a one-loop
$\beta$ function depends on the step convention \SM. The number~100 itself is chosen, not derived
(note in Doc.~133, A040).

\subsection{Geometry, Hilbert space, spectral theory}

The state space is $\mathcal H=L^2(T^4)\otimes\mathbb C^3$; the translation between $T^4$ and
$\mathcal H$ is bijective and lossless (Docs.~230, 270, 330). The torus is flat: the curvature
vanishes pointwise via the metric descent \BM\ (R95). The spectrum separates the lattices:
$\mathbb Z^4$ begins at $|k|^2=1$ with degeneracy~8, $D_4$ at $|k|^2=2$ with 24 (Doc.~314 \KM).

The fundamental fractal operator $\hat F=\bigoplus_n S_n\circ L_n$ with $S_n=r_nU_n$ and
$r_n\in(0,1)$ is bounded \BM. The $\mathbb Z_3$ pairing makes it normal; it is self-adjoint on the
$\mathbb Z_3$-invariant sector $P_0$ \BM, with deficiency indices $(0,0)$ and exactly one
realisation. On the full space only the real part $(\hat F+\hat F^\dagger)/2$ is self-adjoint;
using it is a choice \SM\ (R145, Doc.~327). The particle masses are eigenvalues of the mass
operator $\hat M$, not of $\hat F$ (Docs.~322, 330).

\subsection{Galois structure}

The finite fields $\mathrm{GF}(3^n)$ are the algebraic language of the mode classification on
$T^4/\mathbb Z_3$. One has $|\mathrm{GF}(9)^*|=8$, $|\mathrm{GF}(27)^*|=26=2\cdot13$,
$|\mathrm{GF}(81)^*|=80=2^4\cdot5$; $\mathrm{GF}(9)$ is not a subfield of $\mathrm{GF}(27)$, both
lie in the compositum $\mathrm{GF}(729)$. The Frobenius $x\mapsto x^3$ splits
$\mathrm{GF}(27)^*\cong\mathbb Z_2\times\mathbb Z_{13}$ into two fixed points $\{+1,-1\}$ and
eight three-element orbits; on $\mathbb Z_{13}$ four orbits arise:
\begin{equation}
  O_1=\{1,3,9\},\quad O_2=\{2,5,6\},\quad O_3=\{4,10,12\},\quad O_4=\{7,8,11\}\quad\BM .
\end{equation}
The 5 is the only primitive prime divisor of $3^4-1$, not the smallest one (Docs.~336, 338,
341). With the IPI correspondent D.~Matzke only the algebra is jointly confirmed (orders in $G(3)$
and $G(6)$, Doc.~341); the physical assignments are on the FFGFT side.

% ============================================================
\section{Particles: the ratios}
% ============================================================

\subsection{Lepton ladder}

The charged leptons follow the ladder
\begin{equation}
  \frac{m_i}{v}=r_i\,\xi^{p_i},\qquad r=\Bigl(\tfrac43,\tfrac{16}{5},\tfrac{25}{9}\Bigr),\qquad
  p=\Bigl(\tfrac32,1,\tfrac23\Bigr),\qquad p_g=\tfrac32\bigl(\tfrac23\bigr)^{g-1},
\end{equation}
with $r_g=n_{\phi,g}^2/n_{\theta,g}$ from the winding numbers $(3,2)$, $(5,4)$, $(9,5)$
(Docs.~006, 317, 338) and a geometric sequence of exponents with starting value $D/2$ and ratio
$2/3$ (Docs.~158, 018). The mass ratios contain no scale:
\begin{equation}
  \frac{m_\mu}{m_e}=\frac{12}{5}\xi^{-1/2}=120\sqrt3=207.85\ (+0.52\,\%),\qquad
  \frac{m_\tau}{m_\mu}=\frac{125}{144}\xi^{-1/3}=16.99\ (+1.03\,\%)\quad\KM .
\end{equation}
Formula and ratio calculation are \BM, the agreement \KM; the $r_i$, $p_i$ and winding numbers are
choices \SM\ (R138). The step $1/3$ holds only from $\mu$ to $\tau$; $e$ lies on an orbit of its
own (Doc.~322). The ladder is written without $K_{\text{frak}}$; in ratios it cancels exactly
\BM\ (A040, A100, Docs.~070, 122).

\subsection{Galois identities, Koide and the four routes}

$(m_\mu/m_e)^2_{\text{ladder}}=43200$ is arithmetically exact; against the measured value a
residual of $-1.03\,\%$ remains, an empirical discrepancy \KM\ whose cause is open \SM\ (R113).
The Koide relation, itself a pure ratio,
\begin{equation}
  Q=\frac{m_e+m_\mu+m_\tau}{(\sqrt{m_e}+\sqrt{m_\mu}+\sqrt{m_\tau})^2}=\frac23 ,
\end{equation}
holds for the measured masses to $Q-\tfrac23=-0.017\,\xi$ (PDG 2024; with
$m_\tau=1776.86$\,MeV $-0.046\,\xi$, Doc.~352) \KM; as a prediction of $m_\tau$ from
$m_e,m_\mu$ it lies $0.4\sigma$ from PDG, undecided. From the ladder, Koide follows only
approximately ($+8\,\xi$). In the $\mathbb Z_3$ circulant form
$\sqrt{m_k}=c+d\cos(\theta+2\pi k/3)$ one has $Q=\tfrac23\Leftrightarrow d/c=\sqrt2$ \BM. The
matrix-element route with $|A|=\sqrt2/3$ and $\theta=|A|^2=2/9$ (Doc.~367) hits
$m_\tau/m_e=3477.47$ to $3\times10^{-5}$ \KM\ but gives $m_\mu/m_e=206.7703$, which is
442$\sigma$ from CODATA; the exact pair $(\sqrt2,2/9)$ is excluded for $m_\mu/m_e$ \XM\
(Doc.~369, R121). The geometric derivation of $d/c=\sqrt2$ in Doc.~353 does not hold \SM.
Doc.~369 compares four routes to the lepton masses (ladder, Koide matrix element, Galois,
$\varphi$ skeleton); routes~1 and~3 are the same calculation in two languages for $m_\mu/m_e$.

\subsection{Neutrinos}

Doc.~340 derives hierarchy and mixing from the four Frobenius orbits. The mass ratios are fixed,
\begin{equation}
  m_1:m_2:m_3=1:\sqrt{\tfrac{14}{3}}:11,\qquad
  \frac{\Delta m^2_{\text{atm}}}{\Delta m^2_{\text{sol}}}=\frac{120}{11/3}=32.7\ (-1.4\,\%)\quad\KM ,
\end{equation}
in normal hierarchy. Absolute values require a base mass; the assumed one is
$m_\nu=\tfrac{\xi^2}{2}m_e=4.54$\,meV \SM. This gives $m_2=9.81$ and $m_3=49.96$\,meV,
$\sum m=64.3$\,meV, $\Delta m^2_{\text{atm}}=120\,m_\nu^2=2.476\times10^{-3}$\,eV$^2$ ($-1.0\,\%$
against $2.500\times10^{-3}$, Doc.~340) and $\Delta m^2_{\text{sol}}=\tfrac{11}{3}m_\nu^2=7.565\times10^{-5}$\,eV$^2$
($+0.46\,\%$) \KM. The reactor angle gives $\sin^2\theta_{13}=(25\xi)^{2/3}=0.0223$ ($+1.4\,\%$)
\KM. According to R108, $\nu_1$ and $\nu_2$ lie algebraically in the self-inverse Majorana layer
and $\nu_3$ is Dirac; this assignment of $\nu_3$ to the orbit $O_4$ is in tension with Doc.~346,
which assigns $O_4$ to the charged leptons, and Doc.~344 lists all neutrinos as Dirac \SM; that
the algebraic self-inversion physically means $\nu=\bar\nu$ is a further assumption \SM. For
neutrinoless double beta decay this gives $m_{ee}\approx6.0$\,meV (phases zero), $0.13$\,meV in
the destructive case, at most $14.4$\,meV, and exactly zero if all neutrinos are Dirac
(Doc.~344). Older neutrino spectra (Docs.~006, 007) are excluded \XM\ (R109; note in Doc.~007).
The $A_5$ eight-element spectrum gives $\sin^2\theta_{23}\in\{4/9,5/9\}$; $5/9=0.556$ lies close
to the measured value for normal hierarchy, so maximal mixing $1/2$ is excluded (Doc.~372) \KM.
$\theta_{12}$ is captured only in order of magnitude, the CP phase not at all \SM. The venting
exponent $p=5/24$ follows from the decomposition $\rho_4\!\downarrow_{A_4}=\rho_3\oplus\rho_1$
with $|A_5:A_4|=5$ \BM\ (Doc.~340).

\subsection{Anomalous magnetic moments}

The ratio line carries the result \KM\ (Docs.~018, 158): with winding exponents $5/3$ and $4/3$,
$\Delta a(\tau{-}\mu)/\Delta a(\mu{-}e)=f^{1/3}-1=18.57$, and anchored to the measured $a_e$,
$a_\mu$ one obtains
\begin{equation}
  a_\tau=a_e+\frac{144}{125}\frac{m_\tau}{m_\mu}(a_\mu-a_e)=1.2811\times10^{-3},
\end{equation}
$1.04\times10^{-4}$ above the Standard Model value; present bounds are about 85 times wider, the
decision is pending \SM. A fixed FFGFT addition $\Delta a_\mu=251\times10^{-11}$ as the
explanation of the former muon anomaly is superseded by the 2025 status; it would lie
$3.4\sigma$ above today's difference \XM\ (Doc.~382).

\subsection{Electroweak sector, Higgs, top}

Here too the statements are ratios to the $Z$ mass. On shell, $\sin^2\theta_W=1-M_W^2/M_Z^2=2/9$
holds as an element of the $A_5$ eight-element spectrum (Doc.~372) \KM; this gives
$M_W/M_Z=\sqrt{7/9}$, $M_W=80.42$\,GeV, $3.8\sigma$ from the 2024 world average, an explicitly
stated, falsifiable tension. In the $\overline{\text{MS}}$ scheme the renormalisation running from
$3/8$ gives $\sin^2\theta_W(M_Z)=0.2308$ ($-0.19\,\%$) with two external inputs (Doc.~323) \KM.
The ratios $m_h/M_Z=\tfrac{11}{8}$ ($m_h=125.38$\,GeV, $+0.15\,\%$) and
$m_t/M_Z=(\tfrac{11}{8})^2$ ($m_t=172.40$\,GeV, $-0.10\,\%$) and the trace rule
$M_W^2+M_Z^2+m_h^2=v^2/2$ ($+0.45\,\%$) are candidates \KM; the origin of the factors is a reading
\SM\ (Doc.~373). The chain ``symmetry breaking topologically enforced'' is open at two points
\SM\ (Doc.~355). Earlier Higgs formulas ($m_h=v\xi^{1/4}$, Doc.~041) are refuted \XM\ (R122).

\subsection{Quarks}

The quarks follow the same ladder form $m_q/v=r_q\xi^{p_q}$ with topologically motivated
exponents \KM, except for $c$ and $t$, whose exponents are choices \SM\ (R138); the prefactors
$r_q$ are determined from measured values, their derivation is open \SM\ (R129, Doc.~373). The
reproduction is therefore an encoding, not evidence; the table gives the values with
$v=246$\,GeV:
\begin{center}
\small
\begin{tabular}{@{}lcccc@{}}
\toprule
\textbf{Quark} & $r_q$ & $p_q$ & \textbf{Ladder} ($v=246$\,GeV) & \textbf{PDG 2024} \\
\midrule
$u$ & $6$ & $3/2$ & $2.27$\,MeV & $2.16$\,MeV \\
$d$ & $25/2$ & $3/2$ & $4.73$\,MeV & $4.70$\,MeV \\
$s$ & $26/9$ & $1$ & $94.8$\,MeV & $93.5$\,MeV \\
$c$ & $2$ & $2/3$ & $1.284$\,GeV & $1.273$\,GeV \\
$b$ & $3/2$ & $1/2$ & $4.26$\,GeV & $4.18$\,GeV \\
$t$ & $1/28$ & $-1/3$ & $172.0$\,GeV & $172.6$\,GeV \\
\bottomrule
\end{tabular}
\end{center}
Light quark masses are scheme dependent; $\tilde T\cdot m=1$ refers to pole masses.
$\Lambda_{\text{QCD}}$ does not follow from the $\xi$ ladder \XM\ (R122): $v\,\xi^{1/3}=12.6$\,GeV,
not 200\,MeV. The CKM parameter $\lambda=\xi^{1/6}=0.2260$ lies $0.46\,\%$ from the measured
value \KM; the exponent $1/6$ is a motivated choice \SM\ (note in Doc.~336). The quark and hadron
sector as a whole is an open bridge (Doc.~318, R76).

\subsection{Charge, hypercharge, generations, colour, photon}

Charge quantisation follows algebraically \BM\ (Doc.~346, R111): the QR orbits $O_1$, $O_3$
mod~13 are electrically neutral (neutrinos, antineutrinos), the NQR orbits $O_4$ and $O_2$ are
charged (charged leptons, up quarks); together with $N\bmod3$ this gives
$Q\in\{0,\pm\tfrac13,\pm\tfrac23,\pm1\}$, and ``neutral and coloured'' is excluded. The charge
ratios are thus ratios of integers as well. The down quarks have no orbit of their own.
Gell-Mann--Nishijima is not fully closed: $I_3=\pm\tfrac12$ follows \BM, but the hypercharge of
the charged leptons does not follow from the QR/NQR bit alone and is open \SM\ (R146, Doc.~347).
The three elements of a Frobenius orbit give three generations \BM\ (Doc.~348); the earlier swap
of generations~2 and~3 is an ordering artefact. $\mathfrak{su}(3)$ arises from the
$\mathbb Z_3$ eigensectors \BM; $SU(3)$ is the smallest, not the only suitable group
(Doc.~321). Photon and gluons are massless \BM; that masslessness means spin~1 does not follow
from this and is open \SM\ (Doc.~339).

% ============================================================
\section{The translation into SI}
% ============================================================

\subsection{One anchor}

Because all scales stand in fixed $\xi$ ratios (Section~2), a single measured value suffices for
the translation into SI (Doc.~383). With $m_e$ as anchor one obtains
$E_P=1.237\times10^{19}$\,GeV ($+1.34\,\%$), $G=6.50\times10^{-11}$ ($-2.6\,\%$),
$\ell_P=1.595\times10^{-35}$\,m and $L_0=\xi\ell_P$ \KM. The spread depending on the choice of
anchor is the known residual of the ladder, not a second anchor:
\begin{center}
\small
\begin{tabular}{@{}lccc@{}}
\toprule
\textbf{Anchor} & $E_P$ & $G$ & $\ell_P$, $L_0$ \\
\midrule
$m_e$ & $+1.34\,\%$ & $-2.62\,\%$ & $-1.32\,\%$ \\
$m_\mu$ & $+0.81\,\%$ & $-1.60\,\%$ & $-0.80\,\%$ \\
$m_\tau$ & $-0.22\,\%$ & $+0.45\,\%$ & $+0.22\,\%$ \\
$v$ & $+0.23\,\%$ & $-0.46\,\%$ & $-0.23\,\%$ \\
\bottomrule
\end{tabular}
\end{center}
\noindent This fixes $L_0=2.155\times10^{-39}$\,m, $T_0=7.19\times10^{-48}$\,s and
$E_{\max}=9.16\times10^{22}$\,GeV. With $v=246$\,GeV the absolute lepton masses lie at $-1.18$,
$-0.66$ and $+0.37\,\%$.

\subsection{Electroweak scale}

In the mass calculation $v$ is an intermediate quantity. Depending on the route one obtains:
$248.3$\,GeV with the Galois $m_e$ from Doc.~338 (R104), $248.9$\,GeV with the measured $m_e$,
$245.6$\,GeV with the measured $m_e$ and $K_{\text{frak}}$, $245.65$\,GeV from $\xi^4E_P/(5\pi)$
\KM. The comparison value $v=246.22$\,GeV is a Standard Model definition from $G_F$ \QM\ (R112);
a convention-free comparison would only be possible via the muon lifetime $\tau_\mu$, which is
open \SM.

\subsection{Fractal correction}

\begin{equation}
  K_{\text{frak}}=1-100\,\xi=\frac{74}{75}=0.98\overline{6}.
\end{equation}
The factor appears only in the translation into SI: in ratios it cancels exactly \BM\ (A040,
Docs.~070, 122), and in the basic form $\xi E_0^2=1$ there is no room for it. Its value comes from
the geometry, as the rotation number of the frozen recursion (Section~2). It is confirmed \KM\
(R135): comparing the two routes to $E_0$ with the measured $\alpha$ requires
$K=\xi\,m_em_\mu\,\alpha^{-1}=0.98650$, i.e.\ $n=101.2$ instead of $100$ in $1-n\xi$. This factor
is exactly the square of the bare reference energy in MeV (Section~4.4, Doc.~386, A130). The
residual of $1.7\times10^{-4}$ is not explained by measurement uncertainties and is open \SM. The
\emph{form} of the factor (additive $1-100\xi$ or multiplicative $(1-\xi)^{100}$) is open again
\SM\ (R139): the available witnesses do not separate the two forms more sharply than their
difference of $8.8\times10^{-5}$. Equivalently, $K_{\text{frak}}$ can be written in power form
$(D/3)^{D/2}=74/75$ with $D_f^{\text{eff}}=2.97303$ \KM\ (Doc.~381); the spatial box dimension
$D_f=3-\xi=2.999867$ is a choice \SM\ that does not follow from the 100-step recursion (note in
Doc.~322).

\subsection{\texorpdfstring{$E_0$}{E0} and the SI value of \texorpdfstring{$\alpha$}{alpha}}

In SI, $E_0$ is the geometric mean of the two light lepton masses, bare and with the fractal
correction (R72):
\begin{equation}
  E_0^{\text{bare}}=\sqrt{m_em_\mu}=7.348\ \text{MeV},\qquad
  E_0=\sqrt{\frac{m_em_\mu}{K_{\text{frak}}}}=7.397\ \text{MeV}.
\end{equation}
With $\alpha_{\text{SI}}=\xi\,(E_0/1\,\text{MeV})^2$ one obtains without correction
$\alpha^{-1}=7500/E_0^2=138.9$, with correction $137.06$, $1.7\times10^{-4}$ from the measured
value. The Galois form $\alpha^{-1}=3700/27=137.037$ ($+7.6$\,ppm, Doc.~338, R102) follows from
$(E_0^{\text{bare}})^2=m_em_\mu\approx54\ \text{MeV}^2=|\mathrm{GF}(3)^*|\cdot|\mathrm{GF}(27)|\ \text{MeV}^2$
together with $K_{\text{frak}}=74/75$. Both agreements are properties of the anchor $E_0$, not
independent predictions \KM\ (R135). The SI form requires the reference energy
$E_0\sqrt{\xi/\alpha}=0.99992$\,MeV (bare $0.99323$\,MeV, which is $\sqrt{K_{\text{frak}}}$\,MeV).
The MeV is not a geometric quantity but a unit of the SI chain \QM: its numerical factor is the
numerical value of the historical charge $e$, i.e. of the same unit in which $\alpha$ hides. If the
correction is placed on this side, the side of the Coulomb sphere, the masses and the Compton
sphere stay bare, and $1$\,MeV results to within $8\times10^{-5}$; this assignment is a choice
\SM, the residual remains open (Doc.~386). Third routes to
$E_0$ and closed $\alpha$ formulas from $\xi$ alone are withdrawn \XM.

\subsection{Gravitational constant}

\begin{equation}
  G_{\text{SI}}=\frac{\xi^2}{4m_e}\,C_{\text{conv}}\,K_{\text{frak}},\qquad
  C_{\text{conv}}=7.783\times10^{-3},\quad K_{\text{frak}}=0.986 .
\end{equation}
$G$ is not an independent constant: its form and order of magnitude follow from $\xi$ and the one
anchor \KM\ (Docs.~012, 180, 383). $C_{\text{conv}}$ is the SI conversion of the one-anchor
chain; its corpus value additionally contains the small residual of this chain (chain value
$7.58\times10^{-3}$). The agreement with CODATA to $0.003\,\%$ arises only with this residual and
is not precision evidence in its own right (R141); it holds with the rounded corpus value
$K_{\text{frak}}=0.986$ of the $G$ formula, with $74/75$ it would be $+0.07\,\%$. Explaining the
residual structurally is open \SM. From $\tilde T\cdot m=1$ one has $m\cdot G=\xi^2/4$; a mass
singularity with finite $G$ is excluded \BM\ (Doc.~350), while physical validity in the
black-hole regime is open \SM.

% ============================================================
\section{Gravitation and cosmos}
% ============================================================

\subsection{Black holes}

Equating the Compton and Schwarzschild radii gives, purely geometrically,
\begin{equation}
  M_{\text{coll}}=\frac{m_P}{\sqrt2},\qquad r_s=\sqrt2\,\ell_P,\qquad
  S_{\text{BH}}(M_{\text{coll}})=2\pi\ \text{nat}=\frac{2\pi}{\ln2}\ \text{bit}=9.065\ \text{bit}\quad\BM
\end{equation}
(Docs.~325, 329). $S_{\text{BH}}$ equals Matzke's length in Planck units; Matzke's threshold is
smaller by $\sqrt2$, $n_{\text{thr}}=6.41$ (R144). As a universal bit number $n_{\text{thr}}$ is
refuted \XM\ (R88). The ratio of sector to thermal information is mass independent,
$\ln3\approx1.099$ \BM. The Hawking temperature corresponds to the bit energy $\hbar c/L$ at the
horizon scale \KM; only modes with triality~0 can be emitted, and confinement and Hawking
selection are one principle \BM\ (Doc.~325). Greybody factors and back-reaction are open.

\subsection{Casimir and CMB}

The Casimir energy density between ideal plates is $-\pi^2\hbar c/(720\,d^4)$; the form with 240
is the pressure. At the length $L_\xi$ the ratio to the CMB density, $\pi^2/(720\xi)=102.8$, is
an identity, because $L_\xi$ itself is determined from the CMB density and $\xi$; it is not
measurement evidence (R136). The QFT Casimir energy does not supply a factor $4/3$ \XM\
(Doc.~371). The temperature $k_BT_{\text{CMB}}=\tfrac{16}{9}\xi(1-\tfrac{275}{4}\xi)$\,eV matches
only in the unit eV and with an underived coefficient \SM\ (R137). The factor~3 in front of the
peak sequence $\{1,6,14,26\}$ is a choice; the sequence itself is read off the data (R139).

\subsection{Cosmic sector: deliberately set aside}

The structural form $H_0/\nu_e=\pi^2\xi^{10}=1.75\times10^{-38}$ gives $H_0=66.8$\,km/s/Mpc; the
exponent~10 is a motivated choice (P20), so the statement is conditional and not a prediction.
Equally conditional are $\Lambda^*=1/R_H^2$ and the recasting of the $10^{123}$ question as a
structural question (Doc.~312). Excluded are the $\Lambda$ formulas of older documents,
$H_0=\xi c$, $z=\xi x$ and $z\approx\xi\ln(d/\ell_P)$ \XM\ (R137; note in Doc.~156). The redshift
$1+z=e^{\xi x}$ is achromatic \KM, but the length in $x$ is not specified from $\xi$; the
supernova time dilation is $1+z$ for all $z$ \BM\ (R128). Supernovae, BAO and the CMB
temperature are degenerate between expansion and intrinsic time unfolding; the decision is
theoretical, not empirical (Doc.~267).

\subsection{Galaxies}

The acceleration scale $a_0=c^2\xi^{10}/(4\bar\lambda_e)=1.03\times10^{-10}$\,m/s$^2$ lies
$14\,\%$ below the empirical value and inherits the choice~P20. The gas-dominated dwarf galaxy
DDO\,154 matches to $0.2\,\%$ \KM; for spiral galaxies deviations of about $19\,\%$ remain open.
``Bullet Cluster passed'' is not supported: with baryonic masses the peak lies at the gas (note
in Doc.~308). Light deflection reproduces the GR value $1.75''$ \KM.

% ============================================================
\section{Quantum mechanics}
% ============================================================

FFGFT reads a measurement as a deterministic pole transition in state space,
$A(z,\lambda)=\mathrm{sgn}(z-\lambda)$ with $P(+1)=(1+z)/2$; randomness is epistemic. For several
qubits every such reading is non-local or contextual by Bell and Kochen--Specker, and the $T^4$
reading satisfies this: the correlation is a property of the closed geometry (Doc.~230, R126).
Operationally, FFGFT and standard quantum mechanics give the same predictions up to a $\xi$
correction.

The Bell violation remains: the fractal damping reduces $S$ only by an amount of order $\xi$,
$\Delta S/S\approx10^{-5}$ to $10^{-4}$ depending on the form, and establishes no local bound
(R142). The only documented hardware measurement (IBM, 28~May 2026, Doc.~147, 2048 shots per
run):
\begin{center}
\small
\begin{tabular}{@{}lcccc@{}}
\toprule
\textbf{Device} & \textbf{Runs} & $\langle S\rangle$ & \textbf{Standard error} & \textbf{Share of }$2\sqrt2$ \\
\midrule
ibm\_kingston & 15 & $2.7396$ & $0.0071$ & $96.9\,\%$ \\
ibm\_marrakesh & 10 & $2.7016$ & $0.0138$ & $95.5\,\%$ \\
pooled & 25 & $2.7244$ & $0.0078$ & $96.3\,\%$ \\
\bottomrule
\end{tabular}
\end{center}
\noindent The Bell violation is confirmed \KM; the gap to $2\sqrt2$ ($0.088$ on ibm\_kingston,
$0.104$ pooled) is decoherence, 160 to several thousand times larger than the $\xi$ effect, and
two identical chips differ by $0.038$. The $\xi$ effect cannot be resolved with present-day
hardware \SM. CHSH values given earlier from alleged runs in 2025 and a $\xi$ value fitted to them
are unsupported and removed (R143). The deviation of the $\pi$ gate by $100\xi=1.33\,\%$ is
absorbed by the pulse calibration of real qubits and is not directly measurable (R142). A
factorisation method distinguishable from Shor's does not exist \XM; $\xi$ as a line width of
phase estimation is refuted for the preregistered setup, the line remains open (P44).

% ============================================================
\section{Testable statements}
% ============================================================

\begin{center}
\small
\begin{longtable}{@{}p{3.6cm}p{3.1cm}p{3.6cm}p{2.4cm}@{}}
\toprule
\textbf{Quantity} & \textbf{FFGFT} & \textbf{Measured / bound} & \textbf{Status} \\
\midrule
\endhead
\multicolumn{4}{@{}l}{\textit{Ratios (basic form)}} \\
$m_\mu/m_e$ & $207.85$ & $206.77$ ($+0.52\,\%$) & \KM, $r_i$ \SM \\
$m_\tau/m_\mu$ & $16.99$ & $16.82$ ($+1.03\,\%$) & \KM \\
Koide $Q$ & $2/3$ & $-0.017\,\xi$ & \KM \\
$m_\tau$ from Koide & $1776.97$\,MeV & $1776.93(9)$, $0.4\sigma$ & open \\
$v/E_P$ & $\xi^4/(5\pi)$ & $-0.23\,\%$ & \KM, factor 10 \SM \\
$\xi$ from $m_h/v$ & $1.30\times10^{-4}$ & $4/30000$ ($-2.3\,\%$) & check \KM \\
$\Delta m^2_{\text{atm}}/\Delta m^2_{\text{sol}}$ & $32.7$ & $33.2$ ($-1.4\,\%$) & \KM \\
$\sin^2\theta_{13}$ & $0.0223$ & $0.0220$ ($+1.4\,\%$) & \KM \\
$\sin^2\theta_{23}$ & $5/9$ & normal hierarchy & \KM \\
$a_\tau$ & $1.2811\times10^{-3}$ & bounds 85 times wider & open \\
$M_W/M_Z$ & $\sqrt{7/9}$, $80.42$\,GeV & $3.8\sigma$ (world average) & tension \\
$m_h/M_Z$ & $11/8$, $125.38$\,GeV & $125.20$ ($+0.15\,\%$) & candidate \KM \\
$m_t/M_Z$ & $(11/8)^2$, $172.40$\,GeV & $172.57$ ($-0.10\,\%$) & candidate \KM \\
$\lambda_{\text{CKM}}$ & $\xi^{1/6}=0.2260$ & $0.2250$ ($+0.46\,\%$) & \KM, exponent \SM \\
CHSH deviation & $\Delta S/S\sim10^{-5}$--$10^{-4}$ & hardware noise $\sim10^{-2}$ & not resolvable \\
$\pi$ gate & $1.33\,\%$ & in pulse calibration & not measurable \\[3pt]
\multicolumn{4}{@{}l}{\textit{Absolute values (SI translation)}} \\
$\Delta m^2_{\text{atm}}$ & $2.476\times10^{-3}$\,eV$^2$ & $2.500\times10^{-3}$ ($-1.0\,\%$) & \KM \\
$\Delta m^2_{\text{sol}}$ & $7.565\times10^{-5}$\,eV$^2$ & $+0.46\,\%$ & \KM \\
$m_{ee}$ ($0\nu\beta\beta$) & $6.0$ / $0.13$ / $0$\,meV & nEXO 5--20\,meV & \SM, nEXO decides \\
$\sum m_\nu$ & $64.3$\,meV & cosmological, model dependent & \KM \\
$\alpha^{-1}$ & $137.06$ & $137.036$ ($+1.7\times10^{-4}$) & property of the anchor \\
$\tau_\mu$ without $G_F$, $v$ & -- & MuLan & open \SM \\
\bottomrule
\end{longtable}
\end{center}

% ============================================================
\section{What no longer holds}
% ============================================================

The following statements of older documents are refuted or withdrawn \XM; the documents carry
dated notes on them. In the foundations: $\xi$ as necessarily derived, the shortcut of the Higgs
route to $1/(16\pi^3)$ (the Higgs check to about $2.3\,\%$ remains), a
Casimir factor $4/3$, the number~100 as derived, a minimal time increment $d_1$, self-adjointness
of $\hat F$ on the full space, a field tower $\mathrm{GF}(3)\subset\mathrm{GF}(9)\subset\mathrm{GF}(27)$. For particles: precision tables with fitted
prefactors as evidence, the $\varphi$ mass ladder, the exact pair $(\sqrt2,2/9)$ for $m_\mu/m_e$,
$\Lambda_{\text{QCD}}$ from the $\xi$ ladder, earlier Higgs formulas, the neutrino spectra of
Docs.~006 and 007, a fixed muon addition $\Delta a_\mu=251\times10^{-11}$, Gell-Mann--Nishijima as
fully closed. In gravitation and cosmos: the precision of $G$ as evidence in its own right,
$n_{\text{thr}}$ as a universal number, the Casimir-CMB ratio as measurement evidence, the
$\Lambda$ formulas, $H_0=\xi c$ and $z=\xi x$, ``Bullet Cluster passed''. In quantum mechanics:
local realism or $S\le2+\varepsilon$, a hardware confirmation of the $\xi$ effect, the CHSH values
2.8275 and 2.8278, a factorisation method specific to FFGFT.

% ============================================================
\section{Open bridges}
% ============================================================

Doc.~190 keeps the open points in a list of bridges (as of R146). In summary:
\begin{enumerate}[leftmargin=*,itemsep=1pt]
\item Basic form: the origin of the number~100 and of the factor~10 in $v/E_P$ (Docs.~149, 383);
  $M_{\text{basis}}=m_P$ (R90); self-adjointness of $\hat F$ on the full space (R145); an
  adversarial $D_4$ specificity test (R95); the spectral dimension $d_s=1.86$ versus $\approx2$.
\item SI translation: the residual $1.7\times10^{-4}$ between the $K_{\text{frak}}$ required by
  $\alpha$ and $1-100\xi$, and the form of the factor (R135, R139); the residual of
  $8\times10^{-5}$ in the reference energy $1$\,MeV (Doc.~386); the residual of the one-anchor chain for $G$ (R141).
\item Particles: $\tau_\mu$ directly from FFGFT as the only convention-free test of the weak
  sector (R112); the cause of the decreasing lepton residuals (R113, R121); the Koide amplitude
  $d/c=\sqrt2$ from the geometry; the hypercharge of the charged leptons (R146); the quark
  prefactors and the hadron sector (R76, R129); $K_{\text{had}}$ (R123); two-loop running of
  $\alpha_s$; $m_{\text{Pl}}$ and $\alpha_{\text{em}}(M_Z)$ from $\xi$; spin~1 of the photon.
\item Gravitation and cosmos: greybody factors and back-reaction (R85); a physical Casimir
  signature of $L_\xi$ (R136); $T_{\text{CMB}}$ independent of units (R137); forward selection of
  the CMB peaks and the $C_\ell$ projection (P29--P31); the exponent $41/4$ or 10 forward (P20);
  light elements without a hot early phase; $\Omega_{\text{DM}}$ model-neutrally.
\item Method: $\xi$ as a line width (P44); clean-up of the $H_0$ language (P16, P34).
\end{enumerate}

% ============================================================
\section{Status balance}
% ============================================================

\begin{center}
\small
\begin{longtable}{@{}p{1.2cm}p{13cm}@{}}
\toprule
\textbf{Status} & \textbf{Results (selection)} \\
\midrule
\endhead
\BM & $e^2=4\pi\alpha$ and $\alpha=r_e/\lambda_C$; $\xi E_0^2=1$ in the basic form (Doc.~386);
  charge quantisation from the Legendre symbol (Doc.~346); $SU(2)_L$, $U(1)_Y$ and
  $\sin^2\theta_W|_{\text{GUT}}=3/8$ (R80); $\mathfrak{su}(3)$ from $\mathbb Z_3$ sectors
  (Doc.~321); three generations as a Frobenius orbit (Doc.~348); self-adjointness of $\hat F$ on
  $P_0$ (R145); exponent in $L_0=\xi\ell_P$ (R90); $m_i/E_P=\tfrac{r_i}{5\pi}\xi^{p_i+4}$
  (Doc.~383); rotation number $74/75$ and telescoping product of the recursion (Doc.~381);
  $S_{\text{BH}}(M_{\text{coll}})=2\pi$\,nat and $\ln3$ (R144); supernova time dilation $1+z$
  (R128). \\[3pt]
\KM & Lepton ratios and Koide; neutrino ratios, mass differences and $\theta_{13}$; $a_\tau$
  ratios; $v/E_P$; Higgs check of $\xi$ to about $2.3\,\%$; Higgs and top candidates;
  $\lambda_{\text{CKM}}$; value $K_{\text{frak}}=74/75$ and its appearance only in the SI
  translation; form and order of magnitude of $G$; Bell violation on hardware. \\[3pt]
\SM & $\tilde T\cdot m=1$, $T^4$, $\mathbb Z_3$, $D_4$, $\xi$; $r_i$, $p_i$ and winding numbers;
  the number~100 and the form of $K_{\text{frak}}$; factor~10 in $v/E_P$; $M_{\text{basis}}=m_P$;
  correction on the unit side (Doc.~386); base mass $m_\nu$; quark prefactors; hypercharge of the charged
  leptons; spin~1 of the photon; exponent~10 in the cosmic sector; CMB factor~3. \\[3pt]
\QM & $\hbar=c=1$ and $\alpha=1$ with adjusted $e$; $v=246$\,GeV as Standard Model definition;
  the anchor $m_\mu/m_e$ is QED dependent; Maxwell dynamics adopted; $C_{\text{conv}}$ and the
  MeV as SI conversion. \\[3pt]
\XM & see section ``What no longer holds''. \\
\bottomrule
\end{longtable}
\end{center}

% ============================================================
\section{Method and bookkeeping}
% ============================================================

The corpus labels every statement with a status marker. Newer documents have a verification
script with assertions, many of them null tests. The register Doc.~190 records corrections to
older documents and open bridges; since R133 the source documents themselves are corrected, each
passage with a dated note and the previous value. Excluded predictions are not rewritten but
annotated. From 29~September to 1~October 2026 the entire corpus was recalculated: about 235
documents now carry a correction box, the A series carries notes in the text. The consequences
for the classification are recorded in R135 to R146. Guard rails: a comparison value is not an
input; a form $X=\xi^N$ alone says nothing (P35); predictions are preregistered with fixed
degrees of freedom (P44); every prediction names its anchor (R113); an absolute value in SI tests
the conversion, a ratio tests the structure (A130).

% ============================================================
\section{Reading guide}
% ============================================================

\begin{center}
\small
\begin{tabular}{@{}p{3.2cm}p{11cm}@{}}
\toprule
\textbf{Topic} & \textbf{Documents} \\
\midrule
Foundation & 365 (four layers), 190 (register), 180/189 (chain $\xi\to L_0$), 383 (one anchor) \\
Units, $\alpha$ & A080/A085 (natural units), A130/A267 ($\alpha=1$), 122 (ratio and absolute value), 386 (where $\alpha$ hides), 385 (Higgs--vacuum route) \\
Masses & 006 (ladder), 336/338 (Galois), 351/352 (PDG, lepton residuals), 367/369 (matrix element, four routes) \\
Electroweak & 323 (Weinberg running), 372/373 (2/9, Higgs, Yukawa), 355 (symmetry breaking) \\
Charge, colour & 321 ($SU(3)$), 346/347 (charge, Gell-Mann--Nishijima), 348 (generations) \\
Neutrinos, $g-2$ & 340, 343/344 (Majorana/Dirac), 158/382 ($a_\tau$, status 2025) \\
Gravitation & 012/013 ($G$, SI), 325/329 (black holes), 350 \\
Structure & 314 (lattices), 322/327/330 (spectral theory), 380 ($D_4$), 381 (recursion) \\
Time & 295 (log spiral), 306/307 (reciprocity, time in state space) \\
Quantum mechanics & 230/231 (Hilbert space, measurement), 023 (Bell), 147 (hardware), 175 (gates) \\
Delimitation & 309 (scale anchor), 267 (degeneracy), 377 (interfaces), 378 (peer review) \\
\bottomrule
\end{tabular}
\end{center}

% ============================================================
\section{Conclusion}
% ============================================================

FFGFT is a theory of ratios. In natural units with $\alpha=1$ it replaces the free parameters of
the lepton sector by a structure with one parameter $\xi$ and chosen integers, places lepton
masses, the electroweak scale, the minimal length and the Planck scale in fixed $\xi$ ratios, and
obtains charge quantisation algebraically. The ratios agree to within a per cent down to per
mille. The translation into SI needs one anchor, the fractal correction and conversion quantities;
this is where the absolute values of $\alpha$ and $G$ arise, and where a known residual lies. Open
are the origin of the number~100 and of the factor~10, the form of the fractal correction and its
residual against the MeV, the hypercharge of the charged leptons, the hadron sector and
everything that depends on the cosmic sector. Decisive tests in the coming years are $m_{ee}$
(nEXO), $a_\tau$ (Belle~II), the sum of neutrino masses and $M_W$. The corpus labels choices as
choices and refutations as refutations; this summary reflects that state.

\vspace{1em}
\noindent\textit{Verification script:} \texttt{2/python/Dok384\_Skripte/pruef\_384\_zusammenfassung.py}
--- all numbers in this document recomputed.
